\section{Overlapping normalized windows}

We now pass to bi-infinite sequences by sliding the window along the integers.
This places the construction in the standard framework of factor maps between
shift spaces.

Let $\sigma$ denote the left shift on both $\{0,1\}^{\ZZ}$ and $X_m^{\ZZ}$.

\begin{definition}[Sequence-level fold]
\label{def:Phi-m}
For $m\ge 1$ define
\[
\Phi_m:\{0,1\}^{\ZZ}\to X_m^{\ZZ},
\qquad
(\Phi_m(s))_t:=\Fold_m(s_t s_{t+1}\cdots s_{t+m-1}).
\]
\end{definition}

\begin{proposition}[Finite radius]
\label{prop:Phi-m-radius}
The map $\Phi_m$ is a sliding block code with memory $0$ and anticipation
$m-1$.
\end{proposition}

\begin{proof}
By definition, the symbol $(\Phi_m(s))_t$ depends only on the coordinates
$s_t,\ldots,s_{t+m-1}$.
\end{proof}

\begin{proposition}[Shift equivariance]
\label{prop:Phi-m-equivariant}
For every $m\ge 1$,
\[
\Phi_m\circ \sigma=\sigma\circ \Phi_m.
\]
\end{proposition}

\begin{proof}
Both sides evaluate $\Fold_m$ on the same translated length-$m$ window at each
time $t$.
\end{proof}

\begin{proposition}[Factor property on subshifts]
\label{prop:Phi-m-factor}
Let $S\subset \{0,1\}^{\ZZ}$ be a subshift.  Then $\Phi_m|_S:S\to X_m^{\ZZ}$ is
continuous, shift commuting, and onto its image
\[
Y_m(S):=\Phi_m(S).
\]
Moreover $Y_m(S)$ is a subshift of $X_m^{\ZZ}$, so
\[
\Phi_m|_S:S\twoheadrightarrow Y_m(S)
\]
is a factor map.
\end{proposition}

\begin{proof}
Continuity follows from Proposition~\ref{prop:Phi-m-radius}, and
shift-equivariance from Proposition~\ref{prop:Phi-m-equivariant}.  Since $S$ is
compact and shift invariant, its image under the continuous equivariant map
$\Phi_m$ is again compact and shift invariant.  By definition, the restriction
is onto that image.
\end{proof}

\begin{definition}[The image shift]
\label{def:Y-m}
We write
\[
Y_m:=\Phi_m(\{0,1\}^{\ZZ})\subset X_m^{\ZZ}.
\]
\end{definition}

